\documentclass[../infufrgs/ppgc-tese.tex]{subfiles}
\begin{document}



\section*{Final considerations}
This proposal defines a thesis centered on local interpretability for attention-based models, connecting three fronts in a single empirical program: emergence during training, controllable steering, and diagnostics under shift and perturbation. The core strategy combines sparse feature representations, cross-layer alignment, and logic-based validators so that mechanistic measurements and behavioral outcomes can be analyzed jointly.

Current results already establish feasibility in two directions: controlled delayed-generalization studies under distribution shift, and sparse-feature diagnostics that separate clean from perturbed regimes in vision settings. The remaining work focuses on consolidating the full methodology, executing ablations and baselines across all goals, and reporting reproducible evidence about what is measurable, controllable, and transferable across settings.


\section*{Future work}


The expected outcome is not full global interpretability, but a practical and testable framework for local analysis and intervention, with explicit trade-offs and quantified uncertainty. This scope is aligned with the proposal timeline and with the immediate needs of safer model development and evaluation.

\end{document}
